%==============================================================================
% PREFACE
%==============================================================================
\chapter*{Preface}
\addcontentsline{toc}{chapter}{Preface}

\epigraph{``All models are wrong, but some are useful.''}{--- George Box}

This book is born from a simple observation: the \textit{Notices of the American Mathematical Society} publishes some of the most accessible, beautifully written, and practically relevant expository mathematics anywhere---yet these articles remain scattered across dozens of issues, unknown to many of the scientists and engineers who would benefit most from them.

We have selected fifteen articles published between 2020 and 2026, each chosen for three criteria:
\begin{enumerate}
    \item \textbf{Applied relevance}: The mathematics directly addresses problems in scientific computing, machine learning, engineering, or data science.
    \item \textbf{Implementation readiness}: Each article includes algorithms, code references, or mathematical frameworks that can be translated into working software.
    \item \textbf{Pedagogical clarity}: Written for a broad mathematical audience, these pieces require no specialized background beyond undergraduate linear algebra, calculus, and probability.
\end{enumerate}

The result is a curated textbook that bridges the gap between cutting-edge mathematical research and daily computational practice. Unlike a traditional survey, this book preserves the original authors' voices and narratives---we have lightly edited for consistent notation, added cross-references, and inserted ``Bridge'' sections connecting chapters, but the core exposition remains as published in the \textit{Notices}.

\section*{Organization}

The fifteen chapters are grouped into four parts:

\medskip
\noindent\textbf{Part I: Numerical Foundations} (Chapters 1--5) covers the computational bedrock:
\begin{itemize}
    \item Chapter 1: Why Gaussian elimination---the algorithm behind every linear solver---is both surprisingly stable and dangerously unstable (Urschel, 2025).
    \item Chapter 2: Fewnomial theory for sparse polynomial systems arising in chemical kinetics, power grids, and neural network verification (Rojas, 2025).
    \item Chapter 3: Spectral methods as a unifying language for networks, economics, and opinion dynamics (Golub, 2026).
    \item Chapter 4: 3D printing invariant manifolds of dynamical systems---making the invisible visible (Bishop et al., 2026).
    \item Chapter 5: FAIR principles for mathematical software and digital libraries (Bacher et al., 2026).
\end{itemize}

\medskip
\noindent\textbf{Part II: Machine Learning Theory} (Chapters 6--8) develops a geometric understanding of deep learning:
\begin{itemize}
    \item Chapter 6: Deep networks as affine spline operators---tessellations, loss landscapes, batch norm, and grokking (Balestriero et al., 2025).
    \item Chapter 7: Equivariant neural networks for symmetry-preserving ML in physics, chemistry, and vision (Lim \& Nelson, 2023).
    \item Chapter 8: The Euler Characteristic Transform as a differentiable topological layer (Rieck, 2025).
\end{itemize}

\medskip
\noindent\textbf{Part III: Uncertainty Quantification \& Inverse Problems} (Chapters 9--10) provides complete pipelines:
\begin{itemize}
    \item Chapter 9: UQ tutorial from entropy to EnKF to stochastic surrogate modeling (Chen et al., 2025).
    \item Chapter 10: Tomographic reconstruction of Herculaneum scrolls---a complete inverse problem pipeline (Foschiatti et al., 2025).
\end{itemize}

\medskip
\noindent\textbf{Part IV: Domain Applications} (Chapters 11--15) showcases mathematics in action:
\begin{itemize}
    \item Chapter 11: Algebraic geometry meets GNSS---the global positioning problem as sphere intersection (Boutin \& Kemper, 2025).
    \item Chapter 12: Deep learning for cyber defense---anomaly detection and RL for network security (Emanuello \& Ridley, 2022).
    \item Chapter 13: Coding theory for DNA storage---edit-distance codes, clustering, and reconstruction (Bar-Lev et al., 2022).
    \item Chapter 14: From Hipparchus to Hamiltonian Monte Carlo---inference in astronomy (de Souza et al., 2025).
    \item Chapter 15: Machine-assisted proof---Lean, LLMs, and the future of verification (Tao, 2025).
\end{itemize}

\section*{Companion Resources}

Every chapter has an associated code repository (Python/Julia/MATLAB) at:
\begin{center}
    \url{https://github.com/ams-notices-book/companion-code}
\end{center}
These are not ``supplementary materials''---they are integral to the book. We encourage you to run, modify, and break them.

\section*{Acknowledgments}

We thank the American Mathematical Society for permission to reuse these articles, the authors for writing such clear exposition, and the \textit{Notices} editors (especially Frank Morgan, Karen Saxe, and Daniela De Silva) for curating this body of work. Special thanks to the associate editors who communicated these pieces: Emilie Purvine, Reza Malek-Madani, Vidit Nanda, Richard Levine, Han-Bom Moon, and Laura Schaposnik.

\vspace{1cm}
\noindent\textit{The Editors}\\
\textit{2025}